%% kapa: exam=1 round=34 year=2023 period=1 form=A subject=real_estate subject_ko=부동산학원론 seq=40 number=120
%% source: Q-Net 감정평가사 기출문제 (cst00309) artlSeq=5211789 file=raw/1008/5211789/2023년도 제34회 감정평가사 1차시험 1교시.pdf
%% conversion: kordoc 4.24.51 -> tools/gen1.py (auto); review=proofread-sonnet; flags=none
\begin{question}{120}
감정평가에 관한 규칙의 내용으로 옳지 않은 것은?
\begin{choices}
  \item 시장가치란 감정평가의 대상이 되는 토지등이 통상적인 시장에서 충분한 기간 동안 거래를 위하여 공개된 후 그 대상물건의 내용에 정통한 당사자 사이에 신중하고 자발적인 거래가 있을 경우 성립될 가능성이 가장 높다고 인정되는 대상물건의 가액을 말한다.
  \item 일체로 이용되고 있는 대상물건의 일부분에 대하여 감정평가하여야 할 특수한 목적이나 합리적인 이유가 있는 경우에는 그 부분에 대하여 감정평가할 수 있다.
  \item 감정평가는 대상물건마다 개별로 하여야 하되, 가치를 달리하는 부분은 이를 구분하여 감정평가할 수 있다.
  \item 감정평가법인등은 과수원을 감정평가할 때에 공시지가기준법을 적용해야 한다.
  \item 감정평가는 기준시점에서의 대상물건의 이용상황(불법적이거나 일시적인 이용은 제외한다) 및 공법상 제한을 받는 상태를 기준으로 한다.
\end{choices}
\end{question}
\answer{4}
